% Auto-generated by Code2Latex
% Requires: longtable, booktabs packages

\begin{longtable}{p{\textwidth}}
\caption{Tools for spectra\_elucidation (Level 1)} \\
\endfirsthead

\multicolumn{1}{c}{\tablename\ \thetable{} -- continued from previous page} \\
\endhead

\multicolumn{1}{r}{Continued on next page} \\
\endfoot

\endlastfoot

\textbf{\texttt{carbon\_nmr\_spectra}} \\
\begin{description}
    \item[\textbf{Description:}] Return the 13C NMR spectra for the sample at hand.

This function measures the 13C NMR spectra for the sample at hand. It uses the `get\_c13\_nmr\_prediction` function to measure the experiment NMR spectra. The function returns the 13C NMR spectra as a string.
    \item[\textbf{Arguments:}]
    None
    \item[\textbf{Returns:}] The 13C NMR spectra for the molecule in the sample at hand.

The function returns the 13C NMR spectra as a string. If some error occurs during the NMR spectra generation process, it returns an error message.
\end{description} \\
\midrule
\textbf{\texttt{get\_formula\_from\_smiles}} \\
\begin{description}
    \item[\textbf{Description:}] Generate a chemical formula from a SMILES string using RDKit.

This function takes a SMILES representation of a molecule and generates its chemical formula in Hill notation (C, H, then alphabetical order). It uses RDKit to parse the SMILES string and calculate the molecular formula. If the SMILES string is invalid or cannot be parsed, it returns an error message.
    \item[\textbf{Arguments:}]
    \begin{description}
        \item \texttt{smiles} (str): SMILES representation of a molecule
The SMILES string representing the chemical structure of the molecule. It should be a valid SMILES notation that RDKit can parse.
    \end{description}
    \item[\textbf{Returns:}] The chemical formula in Hill notation (C, H, then alphabetical)

The chemical formula of the molecule represented by the SMILES string, formatted in Hill notation. If the SMILES string is invalid or cannot be parsed, it returns an error message.
\end{description} \\
\midrule
\textbf{\texttt{hsqc\_nmr\_spectra}} \\
\begin{description}
    \item[\textbf{Description:}] Returns the HSQC (Heteronuclear Single Quantum Coherence) NMR spectra for the sample at hand.

This function returns the HSQC NMR spectra for the sample at hand. It returns the HSQC spectrum data as a string, formatted in standard NMR notation, similar to the ACS conventions. If some error occurs during the measurement, it returns an appropriate message.
    \item[\textbf{Arguments:}]
    None
    \item[\textbf{Returns:}] The HSQC NMR spectra for the molecule in the sample at hand.

The function returns the HSQC NMR spectra as a string, formatted in standard NMR notation. If some error occurs during the measurement, it returns an appropriate message.
\end{description} \\
\midrule
\textbf{\texttt{ir\_spectra}} \\
\begin{description}
    \item[\textbf{Description:}] Returns the IR spectra for the sample at hand. This spectra may not be accurate for all compounds. It works best for identifying functional groups such as C=O.

The function returns the IR spectra as a string, following the conventions of IR spectra notation. If some error occurs during the IR spectra measurement process, it returns an error message. Very recommended for C=O group elucidation.
    \item[\textbf{Arguments:}]
    None
    \item[\textbf{Returns:}] The IR spectra for the molecule in the sample at hand.

The function returns the IR spectra as a string, following the conventions of IR spectra notation. If some error occurs during the IR spectra measurement process, it returns an error message.
\end{description} \\
\midrule
\textbf{\texttt{mass\_spectrometry\_spectra}} \\
\begin{description}
    \item[\textbf{Description:}] Returns the mass spectrometry spectra for the sample at hand using the Electrospray Ionization (ESI) technique.

This function returns the mass spectrometry spectra for the sample at hand by making a POST request to an external API that measures the mass spectrometry experiment. The function returns the mass spectrometry spectra as a string in the format "m/z 100.1 (intensity 500), 101.2 (intensity 450), ...".
    \item[\textbf{Arguments:}]
    None
    \item[\textbf{Returns:}] The mass spectrometry spectra for the molecule in the sample at hand.

The function returns the mass spectrometry spectra as a string in the format "m/z 100.1 (intensity 500), 101.2 (intensity 450), ...". If there is some error during the measurement, it returns an appropriate message.
\end{description} \\
\midrule
\textbf{\texttt{obtain\_isomers\_from\_molecular\_formula}} \\
\begin{description}
    \item[\textbf{Description:}] Obtain isomers for a given molecular formula.

This function retrieves isomers for a given molecular formula using the `get\_isomers\_from\_molecular\_formula` remote function. It returns a list of isomer SMILES strings. The list of isomers might not be accurate since it is based on the PubChem database.
    \item[\textbf{Arguments:}]
    \begin{description}
        \item \texttt{molecular\_formula} (str): The molecular formula of the compound for which to retrieve isomers.
The molecular formula representing the chemical composition of the compound for which to retrieve isomers. It should be a valid molecular formula notation that can be processed by the isomer retrieval function.
    \end{description}
    \item[\textbf{Returns:}] A list of SMILES strings representing the isomers of the input compound.

The function returns a list of SMILES strings representing the isomers of the input compound. If no isomers are found, it returns an empty list.
\end{description} \\
\midrule
\textbf{\texttt{proton\_nmr\_spectra}} \\
\begin{description}
    \item[\textbf{Description:}] Returns the 1H NMR spectra for a given SMILES string.

This function returns the 1H NMR spectra for a given SMILES string. It uses the `get\_h\_nmr\_prediction` function to run the experiment NMR spectra. The function returns the 1H NMR spectra as a string.
    \item[\textbf{Arguments:}]
    None
    \item[\textbf{Returns:}] The 1H NMR spectra for the molecule in the sample at hand.

The function returns the 1H NMR spectra as a string, following the conventions of NMR spectra notation. If some error occurs during the NMR spectra generation process, it returns an error message.
\end{description} \\
\midrule
\textbf{\texttt{retrieve\_aromatic\_protons\_shifts}} \\
\begin{description}
    \item[\textbf{Description:}] Retrieve the proton shifts ranges for aromatic hydrocarbons.

This function retrieves the proton chemical shifts ranges for aromatic hydrocarbons. The values represent chemical shift changes (in ppm) caused by substituents on a benzene ring. The values show how much a substituent shifts the resonance of protons at ortho, meta, and para positions relative to unsubstituted benzene. Positive values indicate downfield shifts (deshielding), while negative values indicate upfield shifts (shielding). All shifts are relative to tetramethylsilane (TMS) as the reference standard.
    \item[\textbf{Arguments:}]
    None
    \item[\textbf{Returns:}] A list of dictionaries as a string, containing the substituent effects on proton chemical shifts in aromatic rings

Each dictionary contains the substituent name and its corresponding chemical shift changes (in ppm) for ortho, meta, and para positions. The shifts are based on typical values observed in NMR spectroscopy for various substituents on aromatic rings.
\end{description} \\
\midrule
\textbf{\texttt{retrieve\_carbon\_shifts}} \\
\begin{description}
    \item[\textbf{Description:}] Retrieve the carbon chemical shifts ranges for various functional groups in organic compounds.

This function retrieves the carbon chemical shifts ranges for various functional groups in organic compounds. The chemical shifts (deltas) are reported in parts per million (ppm) relative to tetramethylsilane (TMS) as the reference standard. These values can be used to interpret 13C NMR spectra and identify carbon environments in unknown compounds. It returns a list of dictionaries, each containing the functional group and its corresponding chemical shift range.
    \item[\textbf{Arguments:}]
    None
    \item[\textbf{Returns:}] A list of dictionaries as string containing the carbon chemical shifts ranges for various functional groups in organic compounds

Each dictionary contains the functional group and its corresponding chemical shift range in ppm. The ranges are based on typical values observed in NMR spectroscopy for various functional groups in organic compounds.
\end{description} \\
\midrule
\textbf{\texttt{retrieve\_dbe\_formula}} \\
\begin{description}
    \item[\textbf{Description:}] Retrieve the formula for calculating the Double Bond Equivalent (DBE).

This function returns a string containing the formula for calculating the Double Bond Equivalent (DBE), also known as Degree of Unsaturation (DU) or Index of Hydrogen Deficiency (IHD). The formula is essential in organic chemistry for deducing molecular structures from empirical formulas.
    \item[\textbf{Arguments:}]
    None
    \item[\textbf{Returns:}] A string containing the formula for calculating the Double Bond Equivalent (DBE).

The function returns a string containing the formula for calculating the Double Bond Equivalent (DBE), along with an interpretation of the DBE values and examples of calculations for common organic compounds.
\end{description} \\
\midrule
\textbf{\texttt{retrieve\_isotope\_distribution}} \\
\begin{description}
    \item[\textbf{Description:}] Retrieve the isotopic distribution of common elements in organic chemistry.

This function returns a predefined dictionary containing the isotopic distribution of common elements in organic chemistry, including Carbon, Hydrogen, Sulfur, Chlorine, Bromine, Iodine, Fluorine, Nitrogen, and Oxygen. Each element's isotopes are listed with their natural abundance and m/z values. The function also includes the m/z peaks for each element.
    \item[\textbf{Arguments:}]
    None
    \item[\textbf{Returns:}] A string representation of a dictionary containing the isotopic distribution of common elements in organic chemistry.

The function returns a string representation of a dictionary containing the isotopic distribution of common elements in organic chemistry, including their isotopes, natural abundance, m/z values, and m/z peaks.
\end{description} \\
\midrule
\textbf{\texttt{retrieve\_protons\_shifts}} \\
\begin{description}
    \item[\textbf{Description:}] Retrieve the proton chemical shifts ranges for hydrocarbons.

This function retrieves the proton chemical shifts ranges for various types of hydrocarbons. The chemical shifts (delta) are reported in parts per million (ppm) relative to tetramethylsilane (TMS) as the reference standard. It returns a list of dictionaries, each containing the type of proton and its corresponding chemical shift range.
    \item[\textbf{Arguments:}]
    None
    \item[\textbf{Returns:}] A list of dictionaries as an string containing the proton chemical shifts ranges for hydrocarbons

Each dictionary (as string) contains the type of proton and its corresponding chemical shift range in ppm. The ranges are based on typical values observed in NMR spectroscopy for various types of protons in hydrocarbons.
\end{description} \\
\midrule
\textbf{\texttt{return\_possible\_fragments}} \\
\begin{description}
    \item[\textbf{Description:}] Return a list of fragments of the sample at hand by removing one or two atoms from the molecule.

This function generates some fragments from the sample at hand. It returns a list of unique fragment SMILES strings.
    \item[\textbf{Arguments:}]
    None
    \item[\textbf{Returns:}] A list of unique fragment SMILES strings.

This list contains all the unique SMILES representations of the fragments found by database lookup.
\end{description} \\
\midrule
\textbf{\texttt{search\_by\_smiles}} \\
\begin{description}
    \item[\textbf{Description:}] Search the NMRShift database for entries matching or chemically similar to the given SMILES.

This function searches the NMRShift database for entries that match or are chemically similar to the provided SMILES string. It uses a vector database search to find the top `top\_k` results based on chemical similarity. The function returns a list of matching entries, each containing relevant information about the compound.
    \item[\textbf{Arguments:}]
    \begin{description}
        \item \texttt{smiles} (str): The SMILES representation of the compound to search for in the NMRShift database
The SMILES string representing the chemical structure of the molecule to search for in the NMRShift database. It should be a valid SMILES notation that can be processed by the vector database search.
        \item \texttt{top\_k} (int): The maximum number of results to return. Defaults to 10
The maximum number of search results to return from the NMRShift database. It should be a positive integer.
    \end{description}
    \item[\textbf{Returns:}] A list of dictionaries containing the most relevant entries from the NMRShift database

Each dictionary contains relevant information about the compound, such as its SMILES, chemical shifts, and other properties. The results are sorted by similarity score in descending order.
\end{description} \\
\midrule
\textbf{\texttt{validate\_smiles}} \\
\begin{description}
    \item[\textbf{Description:}] Validate a SMILES string to check if it represents a valid chemical structure.

This function checks if a given SMILES string can be converted into a valid chemical structure using RDKit. It returns True if the SMILES is valid, otherwise returns False.
    \item[\textbf{Arguments:}]
    \begin{description}
        \item \texttt{smiles} (str): The SMILES representation to validate
The SMILES string representing the chemical structure of the molecule to validate. It should be a valid SMILES notation that can be processed by the validation function.
    \end{description}
    \item[\textbf{Returns:}] True if the SMILES string is valid, False otherwise.

The function returns True if the SMILES string can be converted into a valid chemical structure, otherwise it returns False.
\end{description} \\
\end{longtable}
